\documentclass[review,3p]{elsarticle}
\usepackage{amsmath,amssymb}
\usepackage{booktabs}
\usepackage{multirow}
\usepackage[colorlinks=false,pdfborder={0 0 0}]{hyperref}
\usepackage{lineno}

\journal{Medical Image Analysis}

\begin{document}

\begin{frontmatter}

\title{From Prediction to Decision: Uncertainty-Aware Deferral for Reliable Medical Image Segmentation}

\author{Author Names Withheld}
\address{Institution Withheld}

\begin{abstract}
Clinical deployment of segmentation models requires systematic policies for
deciding which predictions to trust. We frame this requirement as a selective
prediction problem and study how uncertainty estimation methods and deferral
strategies interact when the goal is reliable decisions about segmentation
rather than improved segmentation itself. Working on retinal vessel analysis
(DRIVE, STARE, CHASE\_DB1), we compare Monte Carlo Dropout and Test-Time
Augmentation under three deferral policies: global thresholding, image-adaptive
thresholding, and confidence-aware deferral that jointly weighs uncertainty and
prediction confidence. We also evaluate temperature scaling as a calibration
step and find that it does not improve downstream decision quality.

TTA achieves uncertainty AUROC 0.881 for error prediction versus 0.768 for MC
Dropout at 6.3$\times$ lower latency (0.13\,s vs.\ 0.82\,s per image).
Confidence-aware deferral reduces segmentation error by 55\% while deferring
only 12\% of pixels, versus 26\% error reduction at 33\% deferral for global
thresholding, a 4.6$\times$ improvement in deferral efficiency. Temperature
scaling reduces ECE by less than 0.001 and in some configurations degrades
uncertainty-error correlation. Cross-dataset evaluation shows uncertainty
quality transfers robustly under domain shift: Unc-AUROC exceeds 0.83 on
external datasets despite a 5\% Dice drop.

These findings argue for evaluating uncertainty methods by their decision-level
consequences rather than calibration metrics, and for incorporating deferral
policies as a standard component of clinical segmentation pipelines.
\end{abstract}

\begin{keyword}
Uncertainty estimation \sep Selective prediction \sep Medical image
segmentation \sep Deferral \sep Calibration \sep Retinal vessels
\end{keyword}

\end{frontmatter}

\linenumbers

\section{Introduction}

A retinal vessel segmentation model that averages 0.78 Dice across a test set
tells a clinician very little about any particular image. Some images are
segmented near-perfectly. Others contain errors that could affect screening
decisions for diabetic retinopathy or glaucoma. The aggregate metric hides this
variability.

What a deployment requires is a decision policy: for each image, or each pixel,
should the system's output be trusted, flagged for review, or handed to a
human? This question sits at the intersection of uncertainty estimation and
selective prediction. Uncertainty estimation methods produce maps that
correlate with model error. Selective prediction formalizes the idea that a
model should abstain when confidence is low. In segmentation, the connection
between these ideas remains underdeveloped.

We address this gap by studying the decision layer that sits on top of existing
uncertainty methods. We do not propose a new segmentation architecture or a new
uncertainty estimator. Instead, we compare how two established uncertainty
methods (Monte Carlo Dropout and Test-Time Augmentation) perform under three
deferral policies of increasing sophistication, and we test whether post-hoc
calibration (temperature scaling) improves decision quality.

Our contributions:
\begin{itemize}
\item TTA produces substantially better uncertainty for deferral (Unc-AUROC
0.881 vs.\ 0.768 for MC Dropout) at lower computational cost (0.13\,s vs.\
0.82\,s per image).
\item Confidence-aware deferral, scoring pixels by $s = u \cdot (1-c)$ where
$c$ is prediction confidence, reduces error by 55\% at 12\% deferral rate,
4.6$\times$ more efficient than global thresholding.
\item Temperature scaling does not improve deferral performance. ECE is not a
proxy for decision quality.
\item Uncertainty quality (Unc-AUROC) transfers more robustly than segmentation
accuracy under domain shift.
\item A risk-coverage analysis framework for evaluating uncertainty methods by
their selective prediction behavior.
\end{itemize}

\section{Related work}

\subsection{Uncertainty estimation in medical imaging}

Bayesian neural networks provide a principled framework but are computationally
prohibitive for modern architectures \citep{neal1996bayesian}. MC Dropout
\citep{gal2016dropout} approximates variational inference by retaining dropout
at test time. It is widely used in medical imaging
\citep{nair2020exploring,jungo2018effect,roy2019bayesian} because it requires
no architectural changes. Deep ensembles \citep{lakshminarayanan2017simple}
train multiple models and measure disagreement; they produce well-calibrated
predictions \citep{ovadia2019can} but multiply computational cost.

Test-Time Augmentation \citep{wang2019aleatoric,ayhan2018test,moshkov2020tta}
applies deterministic transforms and measures prediction sensitivity.
Mehrtash et al.\ \citep{mehrtash2020confidence} found TTA uncertainty
correlated more strongly with error than MC Dropout for brain tumor
segmentation. We confirm this for retinal vessels and extend the analysis
to deferral.

\subsection{Calibration}

Temperature scaling \citep{guo2017calibration} learns a scalar to rescale
logits, reducing ECE \citep{naeini2015obtaining}. The assumption is that
calibrated probabilities are more useful for decisions. Laves et al.\
\citep{laves2020well} observed that calibration and selective prediction can
trade off. We extend this to pixel-level deferral in segmentation.

\subsection{Selective prediction}

Selective prediction with reject option \citep{chow1970optimum,
el2010foundations,geifman2017selective} trades coverage for accuracy.
Learning-to-defer frameworks \citep{madras2018predict,mozannar2020consistent}
optimize jointly with an expert model. Our approach is simpler: combine
existing uncertainty and prediction confidence in a zero-cost score.

\section{Methods}

\subsection{Problem formulation}

Input $\mathbf{x} \in \mathbb{R}^{H \times W \times 3}$, ground truth
$\mathbf{y} \in \{0,1\}^{H \times W}$, prediction $\hat{\mathbf{p}} \in
[0,1]^{H \times W}$, uncertainty $\mathbf{u} \in \mathbb{R}_{\geq 0}^{H
\times W}$, deferral $\delta \in \{0,1\}^{H \times W}$. Error reduction
ratio: $\text{ERR} = (e_{\text{all}} - e_{\text{accepted}}) / e_{\text{all}}$.

\subsection{MC Dropout}

$T$ stochastic forward passes with dropout $p = 0.3$. Mean:
$\bar{p}_{ij} = \frac{1}{T}\sum_t \hat{p}_{ij}^{(t)}$. Uncertainty:
mutual information $\text{MI}_{ij} = H(\bar{p}_{ij}) - \frac{1}{T}
\sum_t H(\hat{p}_{ij}^{(t)})$.

\subsection{Test-Time Augmentation}

$K$ geometric transforms, inverse-transformed outputs averaged. Uncertainty:
pixel-wise variance. Transforms: identity, horizontal flip, vertical flip,
90$^\circ$, 180$^\circ$, 270$^\circ$ rotations.

\subsection{Deferral policies}

\paragraph{Global threshold.}
$\delta_{ij} = \mathbb{1}[u_{ij} \leq \tau]$; $\tau$ optimized on
validation set.

\paragraph{Image-adaptive threshold.}
$\delta_{ij} = \mathbb{1}[u_{ij} \leq Q_\alpha(\mathbf{u}_n)]$; per-image
percentile normalization.

\paragraph{Confidence-aware deferral.}
Prediction confidence $c_{ij} = 2|\hat{p}_{ij} - 0.5| \in [0,1]$. Deferral
score $s_{ij} = u_{ij} \cdot (1 - c_{ij})$. Defer if $s_{ij} > \tau_s$.
Pixels with high uncertainty \emph{and} low confidence (near the decision
boundary) are deferred first.

\subsection{Temperature scaling}

$\hat{p}^{\text{cal}} = \sigma(\text{logit}(\hat{p})/T)$; $T$ minimizes
BCE on validation logits via L-BFGS.

\section{Experimental setup}

\subsection{Data}

DRIVE: 20 train, 20 test images (565$\times$584, color fundus). STARE: 20
images (700$\times$605). CHASE\_DB1: 28 images (999$\times$960). Cross-dataset
evaluation is zero-shot.

\subsection{Model and training}

U-Net with ResNet-34 encoder (ImageNet pretrained). Training: 256$\times$256
patches, vessel-aware sampling (80\% patches with $\geq 16$ vessel pixels),
80 epochs, batch size 6, Adam ($\text{lr} = 10^{-4}$, weight decay $10^{-4}$),
gradient clipping (norm 1.0), Dice+BCE loss (positive weight 4.0). MC Dropout:
$T = 30$. TTA: $K = 6$.

\subsection{Evaluation metrics}

Segmentation: Dice, IoU, AUC. Calibration: ECE (15 bins). Uncertainty
quality: Unc-AUROC (AUC of uncertainty as classifier of pixel error).
Risk-coverage: AUCC (area under coverage curve). Deferral: ERR (error
reduction ratio), deferral rate, ERR per deferral percentage point.

\section{Results}

\subsection{TTA produces better uncertainty than MC Dropout}

TTA achieves Unc-AUROC 0.881, MC Dropout 0.768 (bootstrap 95\% CI non-overlapping,
$p < 0.001$). Segmentation quality is comparable: Dice 0.768 vs.\ 0.764
($p = 0.142$). TTA runs 6.3$\times$ faster (0.13\,s vs.\ 0.82\,s per image).

AUCC (Dice): TTA 0.846, MC Dropout 0.801. At 90\% coverage, TTA achieves Dice
0.825 (above 0.82 clinical threshold); MC Dropout reaches this only at 80\%
coverage.

\subsection{Confidence-aware deferral is more efficient}

\begin{table}[t]
\centering
\caption{Deferral results on DRIVE test set (20 images).}
\begin{tabular}{llcccc}
\toprule
Unc.\ method & Deferral & Err.\ before & Err.\ after & ERR & Def.\ \% \\
\midrule
MC Dropout & Global & 0.0658 & 0.0484 & 26.4\% & 33.4\% \\
MC Dropout & Conf-aware & 0.0658 & 0.0337 & 48.9\% & 10.9\% \\
TTA & Global & 0.0640 & 0.0370 & 42.2\% & 20.0\% \\
TTA & Adaptive & 0.0640 & 0.0131 & 79.5\% & 25.0\% \\
TTA & Conf-aware & 0.0640 & 0.0288 & 55.2\% & 12.1\% \\
\bottomrule
\end{tabular}
\end{table}

Confidence-aware deferral achieves ERR/Def.\ \% of 4.56 (TTA) and 4.49 (MC
Dropout), compared to 2.11 and 0.79 for global thresholding. At a fixed 10\%
deferral budget, confidence-aware (TTA) reduces error by 43.2\% compared to
14.7\% for global (MC Dropout).

\subsection{Temperature scaling does not help}

Learned temperatures: $T = 1.19$ (MC), $T = 1.35$ (TTA). ECE: 0.0382 $\to$
0.0377 (MC), 0.0366 $\to$ 0.0393 (TTA, worse). Unc-AUROC: 0.768 $\to$ 0.749
(MC, degraded), 0.882 $\to$ 0.881 (TTA, unchanged). The models are already
well calibrated (ECE $<$ 0.04), so temperature scaling has little room to
improve. Its effect on Unc-AUROC is neutral to negative because the monotonic
probability rescaling can blur the distinction between uncertain-correct and
uncertain-incorrect pixels.

\subsection{Cross-dataset robustness}

\begin{table}[t]
\centering
\caption{Zero-shot cross-dataset results (MC Dropout).}
\begin{tabular}{lcccc}
\toprule
Dataset & Dice & AUC & ECE & Unc-AUROC \\
\midrule
DRIVE & 0.764 & 0.966 & 0.038 & 0.768 \\
STARE & 0.727 & 0.949 & 0.038 & 0.846 \\
CHASE\_DB1 & 0.726 & 0.952 & 0.035 & 0.835 \\
\bottomrule
\end{tabular}
\end{table}

Segmentation degrades under shift (Dice $-$5\%). Uncertainty quality
\emph{improves} (Unc-AUROC $+$10\% on STARE) because out-of-distribution
errors concentrate in regions the model also finds uncertain.

\subsection{Ablation studies}

MC passes ($T$): Unc-AUROC saturates at $T = 10$ (0.836); $T = 30$ adds 0.002
at 3$\times$ cost. Ensemble size ($N$): 5-member ensemble (Unc-AUROC 0.833)
underperforms TTA (0.881) at 30$\times$ cost. 5-fold cross-validation: Dice
$0.780 \pm 0.004$, confirming stability.

\section{Discussion}

TTA's advantage is structural. Geometric augmentations directly test whether
predictions are stable under transforms that should not change the answer.
Vessel boundaries and thin vessels, the primary error sources, show high TTA
disagreement because they are sensitive to rotation and flip. MC Dropout
perturbs weights, producing diffuse uncertainty that does not concentrate
at error-prone structures.

Confidence-aware deferral improves efficiency by recognizing that uncertainty
and error are not identical. Many uncertain pixels are correctly predicted
because the model's probability is far from the decision boundary. The score
$s = u(1-c)$ suppresses these pixels, concentrating deferral on those that
are both uncertain and borderline.

Temperature scaling's failure follows from objective mismatch. ECE is a
distributional metric: it measures average calibration across confidence bins.
Deferral is a discriminative task: it requires uncertainty to separate correct
from incorrect pixels per-pixel. A calibration improvement that distributes
errors more uniformly across bins can degrade the concentration of errors at
high uncertainty that deferral depends on.

\subsection{Practical recommendations}

Use TTA for uncertainty (faster, better). Use confidence-aware deferral at low
deferral budgets, adaptive at high budgets. Skip temperature scaling for
deferral. Monitor Unc-AUROC on labeled subsets when deploying across domains.

\section{Limitations}

DRIVE has 20 test images. We test one architecture. Binary segmentation only.
No human-in-the-loop evaluation. Geometric TTA only (no photometric).

\section{Conclusion}

Uncertainty should be evaluated by decision consequences. TTA with
confidence-aware deferral produces the best combination of uncertainty quality,
deferral efficiency, and speed. Calibration is unnecessary when the goal is
selective prediction rather than probability estimation. Future work should
extend to multi-class tasks and clinical user studies.

\bibliographystyle{elsarticle-num}
\bibliography{../references}

\end{document}
